% ===========================================================================
% Figure 1 -- motivating example
%
% Two separate background panels: (a) document, (b) code.
% Three coloured frames, one per inconsistency:
%   1 blue   -- both ends in the document  (inter-document layer)
%   2 orange -- document to code           (document-code layer)
%   3 teal   -- document to code           (document-code layer)
%
% Geometry is driven by explicit coordinates so left/right panels and all
% three frames are perfectly aligned rectangles.
% ===========================================================================
\begin{figure}[t]
  \centering
  \begin{tikzpicture}[
      font=\footnotesize,
      tt/.style={font=\ttfamily\scriptsize, align=left, text=black,
                 inner sep=0pt, outer sep=0pt},
      head/.style={font=\footnotesize\bfseries, text=black,
                   inner sep=0pt, outer sep=0pt},
      lab/.style={font=\scriptsize\bfseries, inner sep=0pt, outer sep=0pt},
    ]

  % -----------------------------------------------------------------------
  % Geometry constants (all in cm)
  % -----------------------------------------------------------------------
  % panels
  % ---- WIDTH KNOBS: shrink these three to narrow the whole figure ----
  \def\PadO{0.17}      % outer padding around each panel
  \def\PadI{0.20}      % inner padding (content inset from panel edge)
  \def\Gap{0.40}       % gap between the two panels
  \def\Lw{6.8}         % left panel content width  (reduces -> text rewraps)
  \def\Rw{5.3}         % right panel content width (reduces -> text rewraps)
  % derived
  \pgfmathsetmacro{\LpanL}{-\PadO}
  \pgfmathsetmacro{\LpanR}{\Lw+\PadO}
  \pgfmathsetmacro{\RpanL}{\Lw+\Gap+\PadO*2+\Gap*0}
  \pgfmathsetmacro{\RpanR}{\RpanL+\Rw+\PadO*2}
  % content x-anchors
  \def\Lx{0}
  \pgfmathsetmacro{\Rx}{\Lw+\Gap+\PadO+\PadI+0.18}

  % vertical anchors (top of each group label)
  \def\Ya{ 0.00}
  \def\Yb{-2.60}
  \def\Yc{-5.78}
  % bottom of last content (set after node placement via coordinate)

  % -----------------------------------------------------------------------
  % Headers
  % -----------------------------------------------------------------------
  \node[head, anchor=south west] at (\Lx, 0.78)
    {(a)\ \texttt{SKILL.md}\,---\,the document};
  \node[head, anchor=south west] at (\Rx, 0.78)
    {(b)\ \texttt{scripts/}\,---\,the code};

  % -----------------------------------------------------------------------
  % GROUP 1  inter-document layer  (blue, left panel only)
  % -----------------------------------------------------------------------
  \node[lab, text=blue!65!black, anchor=north west] (T1) at (\Lx, \Ya)
    {\textcircled{1}\enspace inter-document layer};
  \node[tt, text width=\Lw cm, anchor=north west] (L1)
    at ([yshift=-4pt]T1.south west){%
      :3\enspace "data visualizations with D3.js \colorbox{green!15}{(v6)}."\\[1pt]
      :75\ \ "Pin D3 to a specific version
              (default: \colorbox{red!15}{d3@7.9.0})."\\[1pt]
      \textcolor{black!45}{// one file, two contradictory major versions}};

  % -----------------------------------------------------------------------
  % GROUP 2  document-code layer  (orange)
  % -----------------------------------------------------------------------
  \node[lab, text=orange!80!black, anchor=north west] (T2) at (\Lx, \Yb)
    {\textcircled{2}\enspace document--code layer};
  \node[tt, text width=\Lw cm, anchor=north west] (L2)
    at ([yshift=-4pt]T2.south west){%
      :59\ "Sort input rows deterministically before binding to marks."\\[1pt]
      :60\ "\ldots\,stable grouping order
            (\colorbox{green!15}{Array.from(grouped.keys()).sort()})"\\[1pt]
      \textcolor{black!45}{// the row order must be stable and explicit}};
  \node[tt, text width=\Rw cm, anchor=north west] (R2)
    at ([yshift=-4pt]T2.south west -| \Rx,0){%
      bubble\_chart\_example.js:62\\[1pt]
      \colorbox{red!15}{const sectors =}\\
      \colorbox{red!15}{[...new Set(}\\
      \colorbox{red!15}{data.map(d => d.sector))]}\\[1pt]
      \textcolor{black!45}{// unsorted: order follows input}};

  % -----------------------------------------------------------------------
  % GROUP 3  document-code layer  (teal)
  % -----------------------------------------------------------------------
  \node[lab, text=teal!65!black, anchor=north west] (T3) at (\Lx, \Yc)
    {\textcircled{3}\enspace document--code layer};
  \node[tt, text width=\Lw cm, anchor=north west] (L3)
    at ([yshift=-4pt]T3.south west){%
      :68--71\ "if a force layout is required: fix the
                \colorbox{green!15}{tick count},
                fix initial positions \colorbox{green!15}{deterministically},
                run \colorbox{green!15}{exactly N ticks} and stop."\\[1pt]
      \textcolor{black!45}{// layout output must be reproducible}};
  \node[tt, text width=\Rw cm, anchor=north west] (R3)
    at ([yshift=-4pt]T3.south west -| \Rx,0){%
      bubble\_chart\_example.js:56\\[1pt]
      \colorbox{red!15}{const simulation =}\\
      \colorbox{red!15}{d3.forceSimulation(data)}\\[1pt]
      \textcolor{black!45}{// default tick count, no seed,}\\
      \textcolor{black!45}{// no fixed initial positions}};

  % -----------------------------------------------------------------------
  % Coloured inconsistency frames
  % Each frame spans from its label node to the bottom of its content,
  % left frame: Lx .. LpanR,  cross frame: Lx .. RpanR
  % We derive the bottom from the content nodes.
  % -----------------------------------------------------------------------
  % Frames are fitted to their own content, so the bottom edge follows whichever
  % side is taller.  This stays correct when a different font makes the code
  % column wrap onto more lines (e.g. Latin Modern in Overleaf).
  \node[fit=(T1)(L1), inner sep=4.5pt] (F1) {};
  \draw[blue!60!black, rounded corners=3pt, line width=0.8pt]
    (F1.north west) rectangle (F1.south east);

  \node[fit=(T2)(L2)(R2), inner sep=4.5pt] (F2) {};
  \draw[orange!75!black, rounded corners=3pt, line width=0.8pt]
    (F2.north west) rectangle (F2.south east);

  \node[fit=(T3)(L3)(R3), inner sep=4.5pt] (F3) {};
  \draw[teal!60!black, rounded corners=3pt, line width=0.8pt]
    (F3.north west) rectangle (F3.south east);

  % -----------------------------------------------------------------------
  % Background panels  (drawn last so they go behind frames)
  % We use the bounding boxes of the content nodes plus fixed padding.
  % -----------------------------------------------------------------------
  % Determine overall vertical extent from node positions
  \coordinate (top_y)  at ([yshift=+0.55cm]T1.north);
  \node[fit=(L3)(R3), inner sep=0pt] (BOT) {};
  \coordinate (bot_y)  at ([yshift=-0.30cm]BOT.south);

  \begin{scope}[on background layer]
    % Left panel
    \fill[black!3, rounded corners=4pt]
      (top_y -| \LpanL,0) rectangle (bot_y -| \LpanR,0);
    \draw[black!22, rounded corners=4pt]
      (top_y -| \LpanL,0) rectangle (bot_y -| \LpanR,0);
    % Right panel
    \fill[black!3, rounded corners=4pt]
      (top_y -| \RpanL,0) rectangle (bot_y -| \RpanR,0);
    \draw[black!22, rounded corners=4pt]
      (top_y -| \RpanL,0) rectangle (bot_y -| \RpanR,0);
  \end{scope}

  \end{tikzpicture}

  \caption{Three contract inconsistencies in one real Skill,
    \texttt{d3-visualization}; the document is on the left, the code it ships is
    on the right, and each coloured frame encloses \emph{both ends} of one
    inconsistency, so the frame's horizontal span shows the boundary at which the
    contract breaks.
    \textcircled{1}~\emph{inter-document.} The same \texttt{SKILL.md} names two
    major versions of its own central dependency: line~3 describes the Skill as
    building visualizations ``with D3.js \textbf{(v6)}'', while line~75 pins ``a
    specific version (default: \textbf{d3@7.9.0})'' and line~76 recommends
    vendoring \texttt{vendor/d3.v7.9.0.min.js}. Both ends sit in the document, 72
    lines apart, and nothing else in the Skill says which one holds.
    \textcircled{2}~\emph{document--code.} Lines~59--60 require the reader to
    ``\textbf{sort} input rows deterministically before binding to marks'' and to
    use a ``stable grouping order (explicit
    \texttt{Array.from(grouped.keys()).sort()})''.
    \texttt{bubble\_chart\_example.js} never sorts: it builds the sector list at
    line~62 and the colour-scale domain at line~41 from
    \texttt{[...new Set(data.map(d => d.sector))]}, so both follow the input row
    order. The Skill's \emph{other} example does sort its rows
    (\texttt{interactive\_table\_example.js:73}), so the two shipped examples
    disagree with each other as well.
    \textcircled{3}~\emph{document--code.} Lines~68--71 are conditional and
    binding: ``If a force layout is required: fix the tick count, fix initial
    positions deterministically, run exactly N ticks and stop.''
    \texttt{bubble\_chart\_example.js:56} builds a bare
    \texttt{d3.forceSimulation(data)}. It sets no tick count, no seed and no
    initial positions --- none of the three required actions is taken.
    All three are settled by reading the two files, before any agent runs, and
    none of them raises an error when the Skill is used.}
  \label{fig:motivation}
\end{figure}
